
\documentclass{article}
\usepackage{colortbl}
\usepackage{makecell}
\usepackage{multirow}
\usepackage{supertabular}

\begin{document}

\newcounter{utterance}

\centering \large Interaction Transcript for game `cladder', experiment `full\_v1.5\_default', episode 7019 with Dialogue Architects\_\_qwen3.5{-}9b\_\_SCoRe\_Qwen3.5{-}9B.
\vspace{24pt}

{ \footnotesize  \setcounter{utterance}{1}
\setlength{\tabcolsep}{0pt}
\begin{supertabular}{c@{$\;$}|p{.15\linewidth}@{}p{.15\linewidth}p{.15\linewidth}p{.15\linewidth}p{.15\linewidth}p{.15\linewidth}}
    \# & \multicolumn{2}{c}{Player} && \multicolumn{2}{c}{Game Master} \\
    \hline

    \theutterance \stepcounter{utterance}  
    & & & \multicolumn{4}{p{0.6\linewidth}}{
        \cellcolor[rgb]{0.9,0.9,0.9}{
            \makecell[{{p{\linewidth}}}]{
                \texttt{\tiny{[P1$\langle$GM]}}
                \texttt{You are an expert in causal inference. The following question is not a typical commonsense query, but rather a meticulously designed question created by a professor specializing in causal inference, intended to assess the students' mastery of the course content.} \\
\texttt{Imagine a self{-}contained, hypothetical world with only the following conditions, and without any unmentioned factors or causal relationships: Hwax has a direct effect on gyzp and jyka. Jyka has a direct effect on lirg. Gyzp has a direct effect on lirg. Method 1: We look at how jyka correlates with lirg case by case according to gyzp. Method 2: We look directly at how jyka correlates with lirg in general. To understand how jyka affects lirg, is it more correct to use the Method 1 than Method 2?} \\
\texttt{Start your answer with 'Yes' or 'No', followed by additional reasoning or evidence to support your explanation.} \\
            }
        }
    }
     \\ \\

    \theutterance \stepcounter{utterance}  
    & \multicolumn{4}{p{0.6\linewidth}}{
        \cellcolor[rgb]{0.561,0.737,0.561}{
            \makecell[{{p{\linewidth}}}]{
                \texttt{\tiny{[P1$\rangle$GM]}}
                \texttt{No} \\
\\ 
\texttt{**Reasoning:**} \\
\\ 
\texttt{In this hypothetical world, the causal structure forms a **collider** (or "forking" structure relative to the target variable) configuration regarding the relationship between `jyka` and `lirg`. Specifically, both `jyka` and `gyzp` are direct causes of `lirg`.} \\
\\ 
\texttt{1.  **The Causal Path**: The direct causal effect of `jyka` on `lirg` exists independently of `gyzp`. The path is `jyka → lirg`.} \\
\texttt{2.  **The Role of `gyzp`**: `gyzp` is a confounder in the sense that it is a common cause of `jyka` (via `Hwax → gyzp → jyka`) and `lirg` (via `gyzp → lirg`). However, strictly speaking, `gyzp` is a **confounder** of the relationship between `jyka` and `lirg` because it influences both.} \\
\texttt{3.  **Analysis of Method 2 (General Correlation)**: Method 2 calculates the marginal correlation between `jyka` and `lirg`. Because `gyzp` causes both, there is a spurious association (confounding bias) in the total correlation. The observed correlation will be the sum of the direct effect of `jyka` on `lirg` and the indirect association induced by their shared cause `gyzp`. Therefore, Method 2 does **not** isolate the direct causal effect of `jyka` on `lirg`.} \\
\texttt{4.  **Analysis of Method 1 (Case{-}by{-}Case/Stratification)**: Method 1 looks at the correlation between `jyka` and `lirg` conditional on specific values of `gyzp` (i.e., \$E{[}lirg | jyka, gyzp{]}\$). By conditioning on `gyzp`, we block the non{-}causal backdoor path `jyka ← Hwax → gyzp → lirg`. This isolates the direct causal effect of `jyka` on `lirg`.} \\
\\ 
\texttt{**Conclusion**: To understand *how* `jyka` affects `lirg` (i.e., to estimate the causal effect), one must remove the confounding influence of `gyzp`. Method 2 fails to do this, while Method 1 successfully controls for the confounder. Therefore, Method 1 is the correct approach for causal inference in this scenario.} \\
            }
        }
    }
    & & \\ \\

    \theutterance \stepcounter{utterance}  
    & & & \multicolumn{2}{p{0.3\linewidth}}{
        \cellcolor[rgb]{0.9,0.9,0.9}{
            \makecell[{{p{\linewidth}}}]{
                \texttt{\tiny{[GM$|$GM]}}
                \texttt{No} \\
            }
        }
    }
    & & \\ \\

    \theutterance \stepcounter{utterance}  
    & & & \multicolumn{2}{p{0.3\linewidth}}{
        \cellcolor[rgb]{0.9,0.9,0.9}{
            \makecell[{{p{\linewidth}}}]{
                \texttt{\tiny{[GM$|$GM]}}
                \texttt{yes} \\
            }
        }
    }
    & & \\ \\

    \theutterance \stepcounter{utterance}  
    & & & \multicolumn{2}{p{0.3\linewidth}}{
        \cellcolor[rgb]{0.9,0.9,0.9}{
            \makecell[{{p{\linewidth}}}]{
                \texttt{\tiny{[GM$|$GM]}}
                \texttt{game\_result = LOSE} \\
            }
        }
    }
    & & \\ \\

\end{supertabular}
}

\end{document}
